\section{Analisis dan Pembahasan}
\subsection{Implementasi}

\subsubsection{Modul CNN Layers}
Berikut merupakan kelas-kelas layer yang diimplementasikan from scratch pada modul CNN.

\begin{table}[H]
\centering
\caption{Atribut dan Metode Kelas-kelas CNN}
\begin{tabular}{lp{3cm}p{7cm}}
\hline
\textbf{Kelas} & \textbf{Atribut / Metode} & \textbf{Deskripsi} \\
\hline
\multirow{4}{*}{\code{Conv2D}} & \code{weight, bias} & Parameter bobot konvolusi dan bias.\\
& \code{stride, pad} & Langkah dan padding operasi konvolusi.\\
& \code{\_forward(x)} & Operasi feedforward pada satu buah citra.\\
& \code{forward(x)} & Operasi konvolusi batch / channel support.\\
\hline
\multirow{2}{*}{\code{LocallyConnected2D}} & \code{weight, bias, ...} & Bobot tidak dibagi ukurannya sesuai kernel/patch.\\
& \code{forward(x)} & Forward propagation operasi locally connected.\\
\hline
\multirow{2}{*}{\code{Dense}} & \code{weights, bias} & Bobot linear (fully connected layer).\\
& \code{forward(x)} & Feedforward (X $\cdot$ W + b) dan aktivasi.\\
\hline
\multirow{2}{*}{\code{Flatten}} & \code{forward(x)} & Mengubah tensor multidimensi menjadi vektor 1D.\\
& \code{\_forward(x)} & Implementasi \code{reshape(-1)}.\\
\hline
\multirow{2}{*}{\code{MaxPooling2D}} & \code{pool\_size, stride} & Ukuran dan stride pooling filter.\\
& \code{forward(x)} & Melakukan max pooling window spatial.\\
\hline
\multirow{2}{*}{\code{AveragePooling2D}} & \code{pool\_size, stride} & Ukuran dan stride pooling filter.\\
& \code{forward(x)} & Melakukan average pooling window spatial.\\
\hline
\multirow{1}{*}{\code{GlobalAveragePooling2D}} & \code{forward(x)} & Mean pooling di seluruh dimensi spatial.\\
\hline
\multirow{1}{*}{\code{GlobalMaxPooling2D}} & \code{forward(x)} & Max pooling di seluruh dimensi spatial.\\
\hline
\end{tabular}
\end{table}

\subsubsection{Modul RNN \& LSTM Layers}
Berikut merupakan kelas-kelas layer sekuensial yang diimplementasikan from scratch pada modul RNN dan LSTM.

\begin{table}[H]
\centering
\caption{Atribut dan Metode Kelas-kelas RNN/LSTM}
\begin{tabular}{lp{3cm}p{7cm}}
\hline
\textbf{Kelas} & \textbf{Atribut / Metode} & \textbf{Deskripsi} \\
\hline
\multirow{3}{*}{\code{SimpleRNNCell}} & \code{kernel} & Bobot rekursif dan input konvensional.\\
& \code{recurrent\_kernel} & Bobot internal *hidden state* RNN.\\
& \code{forward(x\_t, h\_prev)} & Perhitungan satu step state RNN.\\
\hline
\multirow{2}{*}{\code{SimpleRNN}} & \code{cell} & Menggunakan \code{SimpleRNNCell}.\\
& \code{forward(x, init\_st)} & Unroll / loop t sequence waktu untuk RNN.\\
\hline
\multirow{3}{*}{\code{LSTMCell}} & \code{kernel, bias} & Bobot gabungan (4 gates) dan bias.\\
& \code{recurrent\_kernel} & Bobot berulang untuk state sebelumnya.\\
& \code{forward(x,h,c)} & Step kalkulasi input, forget, output, control gate. Menghasilkan sel dan *hidden* memori.\\
\hline
\multirow{2}{*}{\code{LSTM}} & \code{cell} & Menggunakan \code{LSTMCell}.\\
& \code{forward(x, init\_st)} & Unroll step-by-step memory tracking.\\
\hline
\multirow{3}{*}{\code{Embedding}} & \code{weights} & Vector vocabulary size $\times$ dimension mapping.\\
& \code{forward(token\_ids)} & Konversi integer token ids menjadi vektor dense.\\
& \code{from\_keras(layer)} & Load transfer matriks embedding Keras.\\
\hline
\end{tabular}
\end{table}

\subsection{Penjelasan Forward Propagation}
% ----------------------------------------------------------
\subsubsection{CNN}
\subsubsection{RNN}
\subsubsection{LSTM}
